\documentclass[10pt,a4paper]{amsart}
\usepackage[margin=19mm]{geometry}
\usepackage{amsmath,amssymb,mathtools}
\usepackage[scr=rsfs,cal=euler]{mathalfa}
\usepackage{xfrac}
\usepackage[unicode,hidelinks,bookmarks=false]{hyperref}
\newcommand{\ie}{i.e.\ }
\newcommand{\cf}{cf.\ }
\newcommand{\resp}{respectively\ }
\newcommand{\Subsection}{\subsection}
\providecommand{\nobreakdash}{\nobreak\hbox{-}}
\setlength{\emergencystretch}{2em}
\multlinegap0pt
\usepackage{amscd}
\usepackage{latexsym,mathrsfs,setspace}
\usepackage{lmodern}
\usepackage{enumitem}
\usepackage{tikz-cd}
\usepackage{setspace}
\usepackage{graphics}
\usepackage{fancyhdr}
\usepackage{tikz}
\usetikzlibrary{decorations.markings,positioning}
\newtheorem{lemma}{Lemma} 
\newtheorem{prop}{Proposition} 
\newtheorem{thm}{Theorem}
\newtheorem{exam}{Example}
\newtheorem{cor}{Corollary}
\newtheorem{defn}{Definition} 
\newcommand{\clc}{{\mathcal C}}
\newcommand{\norm}[1]{\left\lVert#1\right\rVert}
\newtheorem{rmk}[thm]{Remark}
\newtheorem{note}[thm]{Note}
\newtheorem{res}[thm]{Result}
\newcommand*{\tarrow}[2][]{\arrow[Rrightarrow, #1]{#2}\arrow[dash, shorten >= 0.5pt, #1]{#2}}
\newcommand*{\qarrow}[2][]{\arrow[RRightarrow, #1]{#2}\arrow[equal, double equal sign distance, shorten >= 3.3pt, #1]{#2}}
\begin{document}
\title[Categories of bundles and categories of chains]{Categories of bundles and categories of chains}
\author{P G Romeo; Riya Jose}
\date{}
\maketitle
\noindent\emph{Source and licence.} Categories of bundles and categories of chains. Version arXiv:2605.28107v1 (CC BY 4.0). This material is distributed under the Creative Commons Attribution 4.0 International licence, http://creativecommons.org/licenses/by/4.0/, as recorded in the arXiv OAI licence metadata for this exact version and shown on the abstract page https://arxiv.org/abs/2605.28107v1. Full rights notice: licenses/arxiv-2605.28107-CC-BY-4.0.txt.\par\smallskip
\noindent\textit{Editorial scope.} Counted unit: the original \texttt{prop} environment \texttt{-} at source lines 256--258 together with its complete proof at lines 259--283; the delivered document keeps source lines 169--289 so that every referenced label and every citation resolves inside it. Native editable source is the paper's own concatenated TeX; the AMS wrapper is editorial formatting only. No publisher class, upstream script or unrelated artwork is redistributed.
	\section{Category of Bundles}
	In this section we recall concepts of bundles with examples and the category of bundles, for a detailed discussion see \cite{Dale}. We see that category of bundles is a category with subobjects.
	
	\begin{defn}
		A bundle is a triple $(E,p,B)$ where $p: E \rightarrow B$ is a map. The space $B$ is called the base space, the space $E$ is called the total space, and the map $p$ is called the projection of the bundle. For each $b \in B$, the space $p^{-1} (b) $ is called the fibre of the bundle over $b \in B$. 
		\end{defn}
	
	
	\begin{exam}
		The product bundle over $B$ with fibre $F$ is $(B \times F, p,B)$, where $p$ is the projection on the first factor.
	\end{exam}

		A $k-$ dimensional vector bundle $\xi$ over $\textbf{F}$ is a bundle $(E,p,B)$ together with the structure of a $k-$ dimensional vector space over $\textbf{F}$ on each fibre $p^{-1} (b) $ such that the following local triviality condition is satisfied. Each point of $B$ has an open neighbourhood $U$ and a  $U-$ isomorphism $h: U \times F^k \rightarrow p^{-1}(U)$ such that the restriction $ b \times F^k \rightarrow p^{-1} (b)$ is a vector space isomorphism for each $b \in U$.
		
%A space $\textbf{F}$ is the fibre of a bundle  $(E, p, B)$ if every $p^{-1}(b)$ for $b \in B$ is homeomorphic to $\textbf{F}$.  
	\begin{exam}
		The tangent bundle over $S^n$, denoted $\tau (S^n) = (T, p, S^n)$, and the normal bundle over $S^n$, denoted $\nu (S^n) = (N, q, S^n)$, are two subbundles of the product bundle $(S^n \times \mathbb{R}^{n+1}, p, S^n)$ whose total spaces are defined by the relation $(b,x) \in T$ if and only if the inner product $(b|x) = 0$ and by $(b,x) \in N$ if and only if   $x = kb$ for some $k \in \mathbb{R}$. 
		An element $(b,x) \in T$ is called a tangent vector to $S^n$ at $b$, and an  element $(b,x) \in N$ is called a normal vector to $S^n$ at $b$. The fibres $p^{-1}(b) \subset T $ and $q^{-1}(b) \subset N $ are vector spaces of dimension $n$ and $1$, respectively.
	\end{exam}

	
 In the following we recall the definition of a principal $G$- bundle.
	\begin{defn} 
		Let $G$ be a topological group. A principal $G-$ bundle is a map $p: E \rightarrow B$ together with a right action $r: E \times G \rightarrow E$ satisfying the following conditions:\\
		$ 1)\, $ For every $ x \in E $ and $g \in G$, $p(xg) = p(x)$. i.e., $p$ induces a homeomorphism between quotient spaces of this action and $B$.\\
		$2)\, $ For every $b \in B, \, \exists $ an open neighbourhood $U$ and a $G- $ homeomorphism $\phi : p^{-1}(U) \rightarrow U \times G$ (where $G$ acts on $p^{-1}(U)$ by restriction of $r$ on $U \times G$ by $((u,x),g)\longmapsto (u, xg)$) such that following diagram commutes:
		\begin{center}  
			\begin{tikzcd} 
				p{^{-1}(U)} \arrow["\phi "]{r} \arrow["p"]{d} & U \times G  \arrow{dl}{Pr_U}
				\\
				U \\
			\end{tikzcd}		
		\end{center}
	\end{defn}
	\begin{exam}
		Let $B$ be the unit circle $S^1$ parameterized by the angle $\theta$. The total space $E$ is simply the product space $S^1 \times G$, where 
		$G$ is any Lie group. And for simplicity, let's take $G = \mathbb{R}$. The projection map $p : E \rightarrow B$ is defined as $p(\theta, g) = \theta$, where $(\theta, g) \in S^1 \times \mathbb{R}$. The fibre $F$ over each point $b\in B$ is the copy of 
		$G$ at that point.
		In this case, the fibre $F$ over each point 
		$ \theta \in S^1$ is just $G = \mathbb{R}$. The action of the Lie group $G = \mathbb{R}$ on the fibre $F = \mathbb{R}$ is given by addition: $g.f = g + f $, for $g, f \in \mathbb{R}$.\\
		This bundle is called trivial because the total space $E$ is simply the product space $S^1 \times \mathbb{R}$, and each fibre is isomorphic to the structure group $G = \mathbb{R}$. The projection map $p$ is also straightforward, projecting each point to its first component on the base space $S^1$.
	\end{exam}
	\begin{exam}
		The frame bundle of a smooth $n$-dimensional manifold $M$, is a principal $GL(n,\mathbb{R})$-bundle. It consists of all ordered bases (frames) of the tangent spaces of $M$.
		Here the total space is\\ $E = \underset{x\in M}{\cup}\{(x, e_1, e_2, \cdots, e_n)| \{e_1\}_{i=1}^n \textnormal{is a basis of tangent space} \hspace{2mm} T_xM\}$.\\
		Base space is $B = M$. The general linear group \, which represents the set of all changes of basis in $\mathbb{R}^n$ is the fibre. The projection $\pi : E \rightarrow B$ is given by $\pi(x, e_1, e_2, \cdots, e_n)=x$. The general linear group GL(n,R) acts freely and transitively on the fibres. 
	\end{exam}
 


	A space $\textbf{F}$ is the fibre of a bundle  $(E, p, B)$ if every $p^{-1}(b)$ for $b \in B$ is homeomorphic to $\textbf{F}$.  
	Product bundles, tangent bundles, frame bundles discussed above are examples whose fibres are $F$, vector spaces, general linear groups respectively.

	\begin{defn}
		Let $(E, p, B)$ and $(E',p',B')$ be two bundles. A bundle morphism $(u,f) : (E, p, B) \rightarrow (E',p',B')$ is a pair of maps $u: E \rightarrow E'$ and $f: B \rightarrow B'$ such that $up' = pf$.
	\end{defn}
	For two bundles over same base $B$, ie., for  $(E, p, B)$ and $(E',p',B)$ a bundle morphism over $B$ (or $B-$ morphism) $u : (E, p, B) \rightarrow (E',p',B)$ is a map $u: E \rightarrow E'$ such that $up' = p$.	
	In case of morphism of vector bundles the restriction $u: p^{-1}(b) \rightarrow (p')^{-1}(f(b))$ is linear for each $b \in B$ and for $B-$ morphism the restriction $u: p^{-1}(b) \rightarrow (p')^{-1}(b)$ is linear for each $b \in B$.\\
	
%	\begin{prop}
	%	All bundles of the form  $(E, p, B)$ and bundle morphisms forms a category.
%	\end{prop}
	Now, consider a family of bundles $$\{(E, p, B) : E,B \,\text{are total and base spaces, $p$ bundle projection}\}$$
	with the set of all bundle morphisms $(u,f) : (E, p, B) \rightarrow (E',p',B')$ with composition of morphisms  $(u,f) : (E, p, B) \rightarrow (E',p',B')$ and $(u',f') : (E',p',B') \rightarrow (E'',p'',B'')$ 
	 is given by the following commutative diagram: 
	\begin{center}
		\begin{tikzcd}
			E  \arrow["u"]{r}\arrow{d}{p} & E' \arrow["u'"]{r}\arrow{d}{p'} & E'' \arrow{d}{p''}\\
			B  \arrow["f"]{r} & B' \arrow["f'"]{r} & B''\\
		\end{tikzcd}
	\end{center}
	Consequently, composite is the bundle morphism $(uu', ff'):  (E, p, B) \rightarrow (E'',p'',B'')$. The pair $(1_E,1_B):(E, p, B) \rightarrow (E, p, B)$ is a bundle morphism and is the identity morphism. It can be easily seen that all axioms of a category are satisfied.  
	 \begin{prop}
	 	All bundles of the form  $(E, p, B)$ and bundle morphisms forms a category called category of bundles, denoted \textbf{Bun}.
	 \end{prop}
 A bundle morphism $(u,f) : (E, p, B) \rightarrow (E',p',B')$  is an isomorphism if and only if there exists a morphism  $(u',f') : (E',p',B') \rightarrow (E, p, B)$ such that $f'f = 1_B, ff' = 1_B, u'u = 1_E$ and $uu' = 1_E$.
	
	\begin{defn}(cf.\cite{Dale}) \label{subbundle}
	A bundle $(E',p',B')$ is a subbundle of $(E,p,B)$ provided $E'$ is a subspace of $E$, $B'$ is a subspace of $B$, and $p' = p|E' : E' \rightarrow B'$.
\end{defn}
\begin{center}
	\begin{tikzcd}
		E'  \arrow["p'"]{r}\arrow{d}{j_{E'}^{E}} & B'\arrow{d}{j_{B'}^{B}} \\
		E  \arrow["p"]{r} & B \\
	\end{tikzcd}
\end{center}
The bundle morphism $(j_{E'}^{E}, j_{B'}^{B})$ is the inclusion. Being a subbundle is a strict preorder on the vertex class $\nu \textbf{Bun}$.

\begin{prop}
	 The category of bundles \textbf{Bun} with being a subbundle as the choice of subobjects $\mathcal {P}$, is a category with subobjects (\textbf{Bun}, $\mathcal {P}$).
\end{prop}

\begin{proof} Since every bundle is a subbundle of itself and $\mathcal {P}$ being the choice of subbundles we have $\nu \mathcal {P} = \nu \textbf{Bun}$ and morphisms in $\mathcal {P}$  are limited to the canonical inclusions [ see Definition \ref{subbundle}].
	%Let $\mathcal {P}$ be a choice of subbundles in \textbf{Bun}. 
	 Then $(\textbf{Bun},\mathcal {P} )$ is a category with subobjects, since 
	\begin{enumerate}
		\item $\mathcal {P}$ forms a strict preorder with $\nu \mathcal {P} = \nu \textbf{Bun}$.
		\item Every morphism  $(j_{E'}^{E}, j_{B'}^{B}) \in \mathcal {P}$ is a monomorphism being pair of inclusions.
		\item If $ (u, f) = (u'', f'')(u', f') $ where $ (u, f), (u', f') \in  \mathcal {P}$, then $(u'', f'') \in  \mathcal {P}$.\\
		    For, 
		    \begin{center}
		    	\begin{tikzcd}
		    		E''  \arrow["p''"]{r}\arrow{d}{u''} & B''\arrow{d}{f''} \\
		    		E'  \arrow["p'"]{r}\arrow{d}{j_{E'}^{E}} & B'\arrow{d}{j_{B'}^{B}} \\
		    		E  \arrow["p"]{r} & B \\
		    	\end{tikzcd}
		    \end{center}
	    \begin{center}
	    	\begin{tikzcd}
	    		E''  \arrow["p''"]{r}\arrow{d}{j_{E''}^{E}} & B''\arrow{d}{j_{B''}^{B}} \\
	    		E  \arrow["p"]{r} & B \\
	    	\end{tikzcd}
	    \end{center}
    If $ (u, f) = (u'', f'')(u', f') $, then we have $j_{E''}^{E} = u''j_{E'}^{E}$ and $j_{B''}^{B} = f'' j_{B'}^{B}$.\\
    $j_{E''}^{E} = u''j_{E'}^{E} \implies j_{E''}^{E'}j_{E'}^{E} = u''j_{E'}^{E} \implies j_{E''}^{E'} = u''$. Similarly we get, $ j_{B''}^{B'} = f''$. Thus, $(E'',p'',B'')$ becomes a  subbundle of $(E', p', B')$ and hence $(u'', f'') = (j_{E''}^{E'},j_{B''}^{B'} ) \in \mathcal {P}$.
	\end{enumerate}
\end{proof}

\section{Category of chains from bundles}
Consider $\xi_i = E_i \stackrel{p_i}{\rightarrow}  B_i$ of triples $(E_i, p_i, B_i)$ where $p_i$ is a map, $B_i$ is the base space and $E_i$ is the total space. Then the category $\textbf{Bun} $ of bundles is the category whose objects and morphisms are bundles $\xi_i = E_i \stackrel{p_i}{\rightarrow}  B_i$ and bundle morphisms $s_i = (u_i, f_i)$. One remarkable thing regarding the category $\textbf{Bun} $ is that there are many chain  categories associated with them: viz., chains of bundles, chains of fibres and chains of long exact sequences. 
\subsection{Chains of bundles}
%Now we describe chains of bundles consisting of bundles and bundle morphisms between them and it is seen that they form a category. Consider the category $\textbf{Bun} $ of bundles, and a sequence of the form
A chain of bundles \index{chain of bundles} is an indexed family of bundles and bundle morphisms of the form $\cdots \rightarrow (E_i, p_i, B_i) \stackrel{s_i}{\rightarrow} (E_{i+1}, p_{i+1} ,B_{i+1})  \rightarrow \cdots$
%\begin{center}

\begin{thebibliography}{99}
\bibitem[Dale]{Dale} Dale Husemoller, {Fibre Bundles}, {\small \it Springer-Verlag}, {Edition \small \bf 3,}(1994).
\end{thebibliography}
\end{document}
